\subsection{What are the  weak links in Centralized Architecture?}

To systematically identify vulnerabilities within centralized MAS, we leverage our \metricn metric to analyze agent trajectories across different design primitives. Table~\ref{tab:combined_dharma} presents the \metricn classification results across different benchmarks and centralized frameworks, revealing distinct vulnerability patterns that emerge from architectural design decisions.

\textbf{Successful defenses:} We observe that Planner-Stop and Sub-agent Stop occur frequently  across various benchmarks, particularly evident in AgentHarm where 78.74\% of LangGraph cases and 87.40\% of OpenAI Agents cases result in \textbf{Planner-Stop} classification. This occurs because tasks are explicitly harmful and planner successfully recognizes this threat, causing the MAS to halt execution. Sub-agent refusals also serve as effective secondary defenses when planner fail to detect harmful intent, as observed in RedCode's increased \textbf{Sub-agent Stop class}, where coder sub-agent often refuses to generate malicious code despite orchestrators approval. Both of these defenses can be attributed to LLM alignment kicking in, causing the agents to refuse harmful tasks. Below we identify the critical weak links that emerge when these defensive mechanisms fail. 

\begin{wraptable}{r}{0.4\textwidth} % r=right, l=left, width = table width
\centering
\scriptsize
\caption{ARIA scores (\%) for different benchmarks and agent systems with GPT-4o agents. C, D = Centralized, Decentralized.}
\label{tab:eval:results}
\renewcommand{\arraystretch}{1.05}
\setlength{\tabcolsep}{3pt} % tighter column spacing
\begin{tabular}{l|l|rrrr}
\toprule
 & & \multicolumn{4}{c}{ARIA Levels} \\
\cline{3-6}
\rotatebox{90}{} & Framework & 1 & 2 & 3 & 4 \\
\hline
\multirow{6}{*}{\rotatebox{90}{RedCode(exec)}} 
 & SA & 24 & 3.6 & 9.9 & 62.5 \\
 & $C_{\text{Magentic}}$ & 2.6 & 1.0 & 9.3 & 83.70 \\
 & $C_{\text{Langraph}}$ & 8.1 & 3.8 & 34.9 & 52.78 \\
 & $C_{\text{OpenAI}}$ & 6.4 & 0.4 & 23.3 & 69.8 \\
 & $D_{\text{Swarm}}$ & 4.4 & 3 & 19.6 & 73 \\
 & $D_{\text{Langraph}}$ & 5.53 & 29.2 & 40.7 & 24.5 \\
\hline
\multirow{6}{*}{\rotatebox{90}{ASB}} 
 & SA & 6.5 & 0 & 12 & 81.5 \\
 & $C_{\text{Magentic}}$ & 2.99 & 0.25 & 4.24 & 92.52 \\
 & $C_{\text{Langraph}}$ & 23.24 & 0 & 1.13 & 75.63 \\
 & $C_{\text{OpenAI}}$ & 19.6 & 0.3 & 9.4 & 70.7 \\
 & $D_{\text{Swarm}}$ & 0.0 & 0.0 & 1.0 & 99.0 \\
 & $D_{\text{Langraph}}$ & 3.84 & 2.11 & 10.07 & 83.98 \\
\hline
\multirow{6}{*}{\rotatebox{90}{SafeArena}} 
 & SA & 10.4 & 9.6 & 34 & 46 \\
 & $C_{\text{Magentic}}$ & 13.6 & 16.4 & 36.4 & 33.6 \\
 & $C_{\text{Langraph}}$ & 53.2 & 4.8 & 25.2 & 16.8 \\
 & $C_{\text{OpenAI}}$ & - & - & - & - \\
 & $D_{\text{Swarm}}$ & 4.88 & 47.97 & 31.3 & 15.85 \\
 & $D_{\text{Langraph}}$ & 3.6 & 44.4 & 34.4 & 17.6 \\
\hline
\multirow{6}{*}{\rotatebox{90}{AgentHarm}} 
 & SA & 36.4 & 4 & 2.3 & 57.3 \\
 & $C_{\text{Magentic}}$ & 56.2 & 0 & 14.8 & 29 \\
 & $C_{\text{Langraph}}$ & 79.5 & 0 & 1.5 & 19 \\
 & $C_{\text{OpenAI}}$ & 87.4 & 0.8 & 4.7 & 20.5 \\
 & $D_{\text{Swarm}}$ & 1.7 & 22.2 & 10.2 & 65.9 \\
 & $D_{\text{Langraph}}$ & 44.09 & 18.9 & 9.45 & 27.56 \\
\bottomrule
\end{tabular}
\vspace{-5mm}
\end{wraptable}

\textbf{1) Inadequate fallback mechanisms when planning fails enable harmful execution.} The absence of explicit fallback instructions when planner refuses to generate a plan creates a critical vulnerability where systems proceed with task execution despite initial safety rejections. In Magentic, when the planner fails to generate a plan, the system lacks explicit instructions to halt execution and instead invokes subsequent sub-agents in the pipeline. While these sub-agents typically recognize the harmful nature of the task and refuse execution, this behavior results in unnecessary exposure to risk and demonstrates a fundamental design flaw. This pattern is clearly visible in the AgentHarm dataset where Magentic shows increased \textbf{Sub-agent Stop} classification~(51.97\%), indicating that the system relies on downstream agents for safety rather than respecting initial planner refusals. Fig.~\ref{fig:magentic_plan_fail} (Appendix) shows an example agent trajectory from Magnetic Framework exemplifying this behavior.

\textbf{2) Context fragmentation through atomic instruction delegation obscures harmful objectives from sub-agents.} When sub-agents receive only granular, atomic instructions without access to higher-level task context, essentially lacking sub-agent autonomy, they cannot assess the collective harmful intent of their actions. This design choice fundamentally undermines the sub-agents ability to exercise independent safety judgment. In SafeArena benchmark, Magentic's WebSurfer agent receives isolated commands such as "open website" or "click button" without understanding how these actions contribute to a broader harmful objective. This context fragmentation prevents sub-agents from recognizing malicious patterns and results in high Unmitigated-Execution classifications (66.37\%) and increased ARIA-4 risk levels. The architectural decision to maintain orchestrator control through atomic delegation effectively blinds sub-agents to the safety implications of their collective behavior (Fig.~\ref{fig:langgraph_subagent_autonomy} in Appendix).



\textbf{3) Stratified planning architectures enable blind execution of pre-generated harmful plans.} Systems that separate plan generation from execution delegation create vulnerabilities when orchestrators execute plans without reassessing their safety implications. Once a plan is generated and transferred to the orchestrator, the execution phase proceeds mechanically according to the predetermined steps without contextual safety evaluation. This blind adherence to pre-generated plans is particularly problematic in Magentic, where the orchestrator receives a detailed execution plan and delegates tasks based solely on this plan without reconsidering the overall objective's harmfulness. This design flaw contributes to high Unmitigated-Execution classifications (66.37\%) and increased ARIA-4 in SafeArena, as the system loses the opportunity for safety intervention during the  execution phase.

\begin{table}[!h]
\centering
\scriptsize
\caption{\metricn Classification Results for Centralized MAS \small{(GPT-4o agents and GPT-4.1 LLM Judge)}.}
\label{tab:combined_dharma}
\begin{tabular}{@{}ll*{7}{r}@{}}
\toprule
\textbf{Dataset} & \textbf{System} & \multicolumn{3}{c}{\textbf{Planner}} & \multicolumn{2}{c}{\textbf{Sub-agent}} & \textbf{Execution} & \textbf{Errors} \\
\cmidrule(lr){3-5} \cmidrule(lr){6-7}
 &  & \textbf{Stop} & \textbf{Ignore} & \textbf{Fail} & \textbf{Stop} & \textbf{Ignore} & \textbf{} & \textbf{} \\
\midrule
\multirow{3}{*}{ASB} & Magentic & 0.15 & 0.25 & -- & 3.00 & 0.54 & 94.99 & 1.08 \\
 & Langraph & 3.34 & -- & 66.39 & 19.75 & 0.20 & 10.32 & 0.00 \\
 & OpenAI & 14.25 & 43.59 & -- & 4.62 & 0.25 & 35.18 & 2.11 \\
\midrule
\multirow{3}{*}{Agentharm} & Magentic & 3.49 & 0.79 & -- & 51.97 & 0.00 & 26.77 & 16.54 \\
 & Langraph & 78.74 & -- & 0.00 & 4.72 & 0.79 & 15.75 & 0.00 \\
 & OpenAI & 87.40 & 0.00 & -- & 2.36 & 1.57 & 7.87 & 0.79 \\
\midrule
\multirow{3}{*}{RedCode} & Magentic & 3.27 & 3.70 & -- & 8.96 & 4.69 & 79.37 & 0.00 \\
 & Langraph & 4.33 & -- & 0.00 & 13.55 & 1.40 & 80.73 & 0.00 \\
 & OpenAI & 3.30 & 3.70 & -- & 9.00 & 4.70 & 79.40 & 0.00 \\
\midrule
\multirow{3}{*}{SafeArena} & Magentic & 3.14 & 0.00 & -- & 30.49 & -- & 66.37 & 0.00 \\
 & Langraph & 52.40 & -- & 1.20 & 15.60 & 0.40 & 30.40 & 0.00 \\
 & OpenAI & -- & -- & -- & -- & -- & -- & -- \\
\bottomrule
\end{tabular}
\end{table}

\textbf{4)  Sub-agent refusal override represents a critical but infrequent failure mode in orchestrator logic.} While sub-agents successfully refuse harmful requests in most cases, orchestrators occasionally ignore these refusals and proceed with task execution through alternative delegation paths. This failure mode, though less common than the previous weak links, represents a fundamental breakdown in the safety hierarchy of centralized systems. Evidence of this vulnerability appears in RedCode evaluations, where despite clear refusals from Coder or Computer Terminal agents, orchestrators sometimes continue task execution, resulting in "Sub-agent Ignore" score. This indicates insufficient logic for handling and respecting sub-agent safety decisions within the orchestration layer.

\textbf{5) Unmitigated execution reveals fundamental security vulnerability in MAS.} A significant proportion of attack successes manifest as "Unmitigated-Execution" classifications, where no agent within the system recognized or refused the harmful request. The prevalence of this failure mode demonstrates that (i) existing LLM alignment techniques do not reliably transfer to multi-agent contexts; (ii) insufficient context provision to sub-agents, combined with distributed safety responsibility, creates critical security gaps. This enables adversarial prompts to exploit cognitive load distribution, where no single agent maintains sufficient context to identify the overall harmful objective.